\section{CONSIDERAÇÕES FINAIS}
\label{sec:consideracoes}

Este trabalho propôs e aplicou uma abordagem em dois estágios (segmentação socioeconômica por \textit{K-Means} seguida de comparação intragrupo pelo teste de Mann-Whitney) aos 216 municípios maranhenses com dados no IDEB 2025 (anos iniciais do ensino fundamental, rede municipal). A segmentação, feita com a renda \textit{per capita} e a proporção de moradores em pobreza extrema do Censo Demográfico 2022, produziu três grupos (vulnerável, intermediário e de alta renda), partição escolhida tanto pelo \textit{Silhouette Score} quanto pela estabilidade \textit{bootstrap}. Nos dois grupos com amostra suficiente para a análise intragrupo, a \textit{formação docente} (proporção de professores com ensino superior e com licenciatura) está associada ao desempenho de forma consistente. As diferenças entre os municípios de alto e baixo IDEB são de 10 a 12 pontos percentuais, com efeito médio. No grupo intermediário, a presença de quadra esportiva também se destaca. Esses resultados respondem ao objetivo de identificar fatores escolares associados ao desempenho entre municípios de contexto socioeconômico semelhante.

Este estudo apresenta seis limitações relevantes que devem ser consideradas na interpretação dos resultados.
\begin{itemize}
    \item \textbf{Segmentação restrita à dimensão de renda.} As duas variáveis de segmentação descrevem a distribuição de renda, e são correlacionadas entre si ($-0{,}87$). Dimensões socioeconômicas não monetárias, como escolaridade dos adultos e condições de saúde, não entram na segmentação. A ablação mostra que a pobreza extrema acrescenta informação à renda média, mas a partição ainda é essencialmente ordenada ao longo de um eixo de renda.

    \item \textbf{\textit{Cluster} de alta renda fora da análise intragrupo.} Os 15 municípios de alta renda, que reúnem a capital e os principais polos regionais, não têm amostra suficiente para a comparação por quartis. Os fatores associados ao desempenho nesse grupo permanecem sem investigação.

    \item \textbf{Nível de análise.} As variáveis são agregados municipais, e as associações valem para as redes municipais, não para escolas, professores ou alunos individualmente. Além disso, a formação docente da Sinopse inclui os docentes de todas as redes do município; a verificação feita apenas com a rede municipal, que reúne 93\% a 95\% desses docentes nos grupos analisados, reproduz os resultados.

    \item \textbf{Causalidade.} A análise é descritiva e associativa, não causal. Não é possível afirmar que \textit{formação docente} ou \textit{infraestrutura escolar} causam IDEB mais alto. Municípios com gestão mais eficaz tendem a contratar mais docentes qualificados e a obter melhores resultados ao mesmo tempo, tornando difícil separar o efeito de cada variável do efeito da qualidade de gestão municipal.

    \item \textbf{Tamanho das amostras e caráter exploratório.} As comparações envolvem grupos de 23 a 34 municípios. Com a correção para comparações múltiplas, apenas a quadra esportiva e a correção de fluxo no \textit{cluster} intermediário mantêm valor $q$ próximo de 0,05. A formação docente é sustentada pela consistência entre os dois \textit{clusters} e entre os dois procedimentos de comparação, e não por um teste isolado. Os resultados devem ser lidos como exploratórios.

    \item \textbf{Baixo poder explicativo das variáveis disponíveis.} Modelos preditivos treinados com as variáveis do estudo não superaram o acaso na separação dos grupos nem a média na previsão do IDEB. O Censo Escolar não captura processos pedagógicos, gestão, clima escolar ou engajamento familiar, que são mediadores importantes do desempenho e, provavelmente, respondem pela maior parte da variação não explicada.
\end{itemize}

Apesar dessas limitações, a abordagem em dois estágios produziu evidências que podem subsidiar a discussão sobre a alocação de recursos educacionais em diferentes perfis socioeconômicos, ao apontar a formação docente como fator associado ao desempenho de forma recorrente em municípios de renda baixa e intermediária. Verificar se a ampliação da formação docente eleva o desempenho exige estudos com desenho longitudinal ou quase experimental.
